\subsection{Mellomvare, REST og adapter-mønsteret}
\label{subsec:mellomvare-rest-adapter}

Mellomvare er programvare som fungerer som et kommunikasjonslag mellom applikasjoner,
databaser og operativsystemer. I stedet for at hver applikasjon må koble seg direkte til
alle andre, kobler de seg til mellomvaren, som håndterer ruting, transformasjon og
sikkerhet på tvers av systemene \cite{azure_middleware}.

\gls{rest} er en arkitekturstil for distribuerte systemer,
definert av Roy Fielding i hans doktoravhandling fra 2000. REST bygger på et sett
constraints: klient-tjener-separasjon, tilstandsløshet, og et uniformt grensesnitt
basert på ressursidentifikasjon via URI-er \cite{wiki_rest}. VideoAPI følger disse
prinsippene ved å eksponere videorom som ressurser over HTTP.

Adapter-mønsteret er ett av de klassiske strukturelle designmønstrene beskrevet av
Gamma et al. Det løser problemet der to grensesnitt er inkompatible ved å innføre
et mellomliggende adapter-lag som oversetter kall fra klienten til det formatet
mottakeren forventer \cite{wiki_adapter}. I VideoAPI realiseres dette gjennom
\texttt{IVideoProvider}-grensesnittet, der hver provider-implementasjon fungerer som
en adapter mot sin respektive videomotor.

\subsection{Containerisering og statelessness}
\label{subsec:containerisering-statelessness}

Containerisering innebærer at en applikasjon pakkes sammen med alle sine avhengigheter
i et isolert miljø som kan kjøres konsistent på tvers av ulike plattformer \cite{docker_docs}.
Docker er den mest utbredte teknologien for dette, og muliggjør reproduserbare bygge-
og kjøremiljøer uten avhengighet av vertssystemets konfigurasjon.

Twelve-Factor App-metodologien definerer et sett prinsipper for å bygge skalerbare
og vedlikeholdbare tjenester. To av prinsippene er særlig relevante for VideoAPI:
konfigurasjon skal lagres i miljøvariabler fremfor i kode, og tjenesten skal være
tilstandsløs slik at flere instanser kan kjøre parallelt uten å dele \gls{prosess-intern-state} \cite{twelve_factor}. VideoAPI følger begge disse prinsippene — all konfigurasjon
styres via miljøvariabler, og API-et behandler hver forespørsel uavhengig av tidligere
forespørsler.

\subsection{Observability med Prometheus og Grafana}
\label{subsec:observability}

Observability handler om evnen til å forstå tilstanden til et system basert på
data det eksponerer. Prometheus er et åpen kildekode-verktøy for systemovervåkning
som samler inn og lagrer metrikker som tidsseriedata, identifisert med metrikknavn
og nøkkel-verdi-par kalt labels \cite{prometheus_docs}. Innsamling skjer via en
pull-modell der Prometheus periodisk henter metrikker fra instrumenterte tjenester.

Grafana brukes til å visualisere disse metrikkene gjennom dashboards sammensatt av
paneler som spør mot én eller flere datakilder \cite{grafana_docs}. I VideoAPI
eksponerer applikasjonen metrikker via \texttt{prometheus-net}, som Prometheus
skraper og Grafana visualiserer i egne dashboards for drifts- og sikkerhetsovervåkning.

\subsection{Tokenbasert autentisering}
\label{subsec:tokenbasert-autentisering}

OAuth 2.0 er et autorisasjonsrammeverk som lar en klient få begrenset tilgang til
en tjeneste på vegne av en ressurseier \cite{rfc6749}. I maskin-til-maskin-sammenheng
benyttes \texttt{client\_credentials}-flyten, der klienten autentiserer seg direkte
mot en autorisasjonsserver og mottar et tilgangstoken uten brukerinteraksjon.

JSON Web Token (JWT) er et kompakt, URL-trygt format for å overføre påstander mellom
parter som et signert objekt \cite{rfc7519}. Tokenet består av tre deler: header,
payload og signatur. I VideoAPI valideres JWT-tokens utstedt av Keycloak på hvert
beskyttet endepunkt, og tilgang krever at tokenet inneholder \texttt{api\_access}-påstanden.